\chapter{Principios Forenses y Cadena de Custodia}

La validez jurídica y técnica de un dictamen forense descansa en la inviolabilidad de la evidencia. Frente a una sospecha de compromiso, el perito debe resolver el dilema fundamental entre preservar el estado volátil de la máquina en funcionamiento o proceder al apagado controlado para su examen estático en laboratorio.

\begin{notaforense}[La Regla de Oro Forense]
Nunca se investiga directamente sobre el medio original de almacenamiento si no es estrictamente imprescindible. La clonación bit a bit y el cotejo criptográfico mediante sumas de verificación MD5 y SHA1 garantizan que el estado de la evidencia permanezca inalterado a lo largo del proceso pericial.
\end{notaforense}

\section{Análisis \emph{post-mortem}}


 Desde que se sembraron los programas de registro se hizo una copia
completa del sistema, y antes del análisis forense se generaron varias
copias idénticas bit a bit de los discos duros. Salvo por necesidad
se trabajo en el disco original si era meritorio. La copia de las
particiones del disco duro se analizó en sistemas limpios y
controlados.

\subsection{Sistemas de análisis y referencia}


\textbf{Dell precision 350}\\
Pentium 4 2.8 GHz\\
512 MB RAM\\
DD 80 GB\\
Debian GNU/Linux Woody (referencia.)\\
Debian GNU/linux Sarge (análisis.)\\

\textbf{Dell inspiron 4150}\\
Pentium 4 1.8 GHz\\
256 MB RAM\\
DD 30 GB\\
Debian GNU/Linux Woody (referencia.)\\
Debian GNU/linux Sarge (análisis.)\\


\textbf{Programas usados para el análisis}\\



\vspace{1em}
\vspace{0.8em}\hrulefill\vspace{0.8em}
{\footnotesize
\begin{lstlisting}[style=terminal]

-- strings 
								  
-- strace 

-- ltrace

-- gdb

-- find 

-- md5sum

-- netstat 

-- /usr/sbin/lsof 

-- grep 

-- dd 

-- export 

-- chkrootkit 

-- icat 

-- file 
    
-- readelf 

-- objdump 
   
-- ps 
    
-- nc 
    
-- stat 

-- nm 

-- unrm 

-- mactime 

-- sleuthkit

-- autopsia

-- entropy

-- gcc

-- f-prot 

-- lsattr

-- lazarus

-- dcat                            

-- dls                              

-- dstat                            

-- ffind                            

-- file                             

-- fls                              

-- fsstat                           

-- hfind                            

-- icat                             

-- ifind                            

-- ils                              

-- istat                            

-- mactime                          

-- md5                              

-- sha1                             

-- sorter

-- grave-robber

-- lastcomm

-- major_minor

-- timeout

-- bdf                                                      

-- ils2mac  
                                          
-- realpath

-- check_wtmpx                                          

-- chkdirs                                                   

-- chklastlog                                                

-- chkpro                                                  

-- chkproc                                                  

-- chkwtmp                                                  

-- ifpromisc                                                

-- snif                                                     

-- last

\end{lstlisting}
}
\vspace{0.8em}\hrulefill\vspace{0.8em}

\begin{perspectiva20}[Perspectiva 20 Años Después: Adquisición de Memoria Volátil]
En 2004, la adquisición de memoria RAM viva estaba en una fase incipiente (herramientas como \textit{fmem} o \textit{dd} sobre \texttt{/dev/mem}). Hoy en día, la captura de memoria volátil mediante herramientas como \textit{LiME} (\textit{Linux Memory Extractor}) y su posterior análisis con \textit{Volatility 3} constituyen el primer paso ineludible antes de apagar cualquier sistema comprometido, permitiendo rescatar llaves criptográficas, sockets y procesos inyectados que desaparecen al apagar el equipo.
\end{perspectiva20}
